% Source presentations for register row N-c-014.

\subsection{Finite-word monadic logic and Rabin's infinite-tree theorem}
\label{sec:p09-presentation-n-c-014}

\subsubsection{The three original predicate-and-automata presentations}

The registered item is B\"uchi's 1960 journal article, but the complete body
available for reconstruction is its September 1959 University of Michigan
technical-report precursor, UMR1083, under the same title.  The report has 45
numbered pages and contains the complete development from syntax through the
finite-word theorem and bibliography.  The discussion below therefore cites
the report's own page numbers and theorem numbers.  It does not claim that the
report and the 1960 article have identical wording, numbering, examples, or
proof detail; that page-for-page journal reconciliation remains open
\cite{Buchi1959WeakReport,Buchi1960}.

\paragraph{The source's weak second-order arithmetic.}
B\"uchi begins with an interpreted monadic second-order language.  Individual
variables range over \(\mathbb N\); \(0\) and successor are primitive; and a
predicate variable ranges over what the source calls a \emph{special
predicate}: a Boolean-valued function on \(\mathbb N\) which is eventually
false, equivalently a finite subset of \(\mathbb N\).  An \(n\)-predicate is
an \(n\)-tuple of such predicates.  The interpreted system is denoted
\(W.2.A.\); B\"uchi expressly supplies semantics rather than a deductive
calculus.  Thus \emph{weak} here means finite-set second-order quantification.
It does not mean a restriction on quantifier alternation, and it is not the
full-set infinite-word formalism of his later paper
\cite[pp.~2--6]{Buchi1959WeakReport}.

\paragraph{The source automaton and its behavior.}
For input truth-value variables \(X^n\) and transit-state variables \(Y^m\),
an \(n/m\)-automaton is the pair of propositional-formula tuples
\[
 E^m,
 \qquad
 H^m[Y^m,X^n].
\]
Here \(E^m\) fixes the initial \(m\)-bit state and \(H^m\) gives the next
\(m\)-bit state.  If \(i\) is a special \(n\)-predicate, its transit predicate
\(r\) is the unique recursion
\[
 r(0)=E,
 \qquad
 r(t+1)=H[r(t),i(t)].
\]
A binary output is a propositional formula \(U[Y^m]\), and the corresponding
output predicate is \(u(t)=U[r(t)]\).  The output is \emph{special} for
\((E,H)\) when the propositional identity
\[
 U[Y^m]
 \equiv
 U[H(Y^m,\bot^n)]
\]
holds identically in the state variables \(Y^m\).  This says that one
all-false-input transition preserves the output value at every state.  Once a
special input has become permanently false, the identity makes \(u\)
eventually constant.
The behavior is then
\[
 \beta(E,H,U)
 =
 \left\{
   i:
   i\text{ is special and }
   \exists x\,\forall t\geq x\;U[r(t)]
 \right\}.
\]
Thus \(E,H,U\) are finite propositional data, \(i,r,u\) are predicates on
time, and \(\beta(E,H,U)\) is a set of finite-set assignments; none of these
standings is interchangeable \cite[pp.~6--8]{Buchi1959WeakReport}.

\paragraph{Synthesis and analysis.}
Lemma 1 effectively converts every \(W.2.A.\) formula \(C(i)\) to a local
successor form
\[
 (\mathbf Q\mathbf j)\,
 \left[
   K(\mathbf j(0))
   \land
   \forall t\,
   B(i(t),\mathbf j(t),\mathbf j(t+1))
 \right],
\]
where \(\mathbf Q\mathbf j\) is a finite prefix of finite-set quantifiers and
\(K,B\) are propositional matrices.  Theorem 1 then gives the effective
\emph{synthesis} direction: from any \(C(i)\), construct a finite automaton
with special output whose behavior is exactly the special assignments
satisfying \(C\).  The proof first handles the quantifier-free local matrix,
uses output negation for complement, replaces arbitrary propositional
transition/output data by equivalent disjunctive state formulae, and removes
one input component for an existential finite-set quantifier.  Alternating
these operations realizes the quantifier prefix.

The effective \emph{analysis} direction constructs, from an
\(n/m\)-automaton with special output, a formula in the stronger normal form
\[
 C(i)
 =
 (\exists\mathbf j)\,
 \left[
   K(\mathbf j(0))
   \land
   \forall t\,
   B(i(t),\mathbf j(t),\mathbf j(t+1))
 \right]
\]
with \(i\models C\) exactly when \(i\in\beta(E,H,U)\).
Quantify finite predicates \(r,j\), require \(j(0)\), and use
\[
\begin{aligned}
 r(0)&=E,\\
 \forall t\;&\bigl(j(t+1)\rightarrow j(t)\bigr),\\
 \forall t\;&\bigl(j(t+1)\rightarrow
     r(t+1)=H(r(t),i(t))\bigr),\\
 \forall t\;&\bigl(\neg j(t+1)\rightarrow i(t)=\bot^n\bigr),\\
 \forall t\;&\bigl(j(t)\land\neg j(t+1)\rightarrow U(r(t))\bigr).
\end{aligned}
\]
Because \(j\) is finite, nonempty, and downward closed, it has a last true
point \(h\).  The finite witness \(r\) follows the actual run only through
\(h\); the input is false from \(h\) onward; and the boundary clause together
with the pointwise special-output identity propagates truth along the omitted
tail.  Conversely, for an accepted special input choose \(h\) after both input
extinction and eventual output truth, retain the actual run through \(h\), and
set the witness false afterward.  These local clauses therefore prove the
analysis direction of Theorem~2
\cite[pp.~14--18, thms.~1--2]{Buchi1959WeakReport}.

\paragraph{The source's finite-word encoding.}
For an integer \(n\geq1\), let \(\Sigma_n=\{\bot,\top\}^n\).  Section 7 maps a nonempty finite word
\(Y_0\cdots Y_h\in\Sigma_n^*\) to the special \(n\)-predicate
\[
 f(Y_0\cdots Y_h)(t)
 =
 \begin{cases}
   Y_t,    &0\leq t\leq h,\\
   \top^n, &t=h+1,\\
   \bot^n, &t>h+1.
 \end{cases}
\]
The empty word receives the terminal marker \(\top^n\) at \(0\) and is false
afterward.  The marker is part of the encoding, not an unmentioned end position.
For an arbitrary output \(U\), the source's \emph{string behavior} accepts a
word \(w\in\Sigma_n^*\) exactly when \(U[r(|w|)]\) holds for the run \(r\) on
its encoded input.  Thus a nonempty \(w=Y_0\cdots Y_h\) is read at
\(r(h+1)\), whereas the empty word is read at \(r(0)\).  Theorem 6 proves that a set is
such a string behavior exactly when it is the behavior of an automaton with
special output.  Theorem 7 then gives effective formula-to-automaton and
automaton-to-formula translations for these encoded finite words.  Combining
that result with the already established automata characterization of regular
sets, Theorem 8 makes the four conditions equivalent: arbitrary-output string
behavior, special-output behavior, \(W.2.A.\)-definability by a formula whose
assignments are word encodings, and regularity
\cite[pp.~30--34, thms.~6--8]{Buchi1959WeakReport}.
The analysis half of Theorem~7 uses the finite-witness construction above.

\paragraph{The decision proof and arithmetic consequences.}
Theorem 3 extends the local normal form to formulas with free individual
variables.  Represent \(x\) by a downward-closed special predicate satisfying
\(s(x)\land\neg s(x+1)\).  Its support is exactly \(\{0,\ldots,x\}\),
including the singleton \(\{0\}\) when \(x=0\).

Corollary 1 decides closed \(W.2.A.\) sentences.  For a normal sentence
\(\exists\mathbf j[K(\mathbf j(0))\land\forall t\,
B(\mathbf j(t),\mathbf j(t+1))]\), first check
\(B(\bot,\bot)\).  If it is false, no finite predicate can satisfy the
infinite all-false tail.  If it is true, let \(k\) be the number of Boolean
states and search for \(x\leq k\) and a path
\(Y_0,\ldots,Y_x\) which starts in a \(K\)-state, follows \(B\), and ends at
\(\bot\).  A path with repetition can have the intervening loop deleted, so
this is a finite search.  The false-tail guard and the source's inclusive bound
are both essential parts of the decision proof
\cite[pp.~18--20]{Buchi1959WeakReport}.

Corollary 2 proves that a unary subset of \(\mathbb N\) is definable in
\(W.2.A.\) exactly when it is ultimately periodic, and its proof effectively
extracts a phase and period.

\paragraph{Binary arithmetic and finite-set arithmetic.}
Encode a finite predicate by the bijection
\[
 F(i)=\sum_{t\in i}2^t.
\]
Use the arithmetic structure \(\mathcal B_2=(\mathbb N,+,\in_2)\), where
\[
 x\mathrel{\in_2}y
 \quad\Longleftrightarrow\quad
 \exists t\in\mathbb N\;
   \bigl(x=2^t\text{ and the coefficient of }2^t\text{ in }y\text{ is }1\bigr).
\]
This binary relation is primitive; the displayed powers describe its
interpretation.  The predicate \(\operatorname{Pow}(x):=(x\in_2x)\)
defines exactly the positive powers of two.  The constants \(0,1\) and
order used below are definable from addition: zero satisfies \(z+z=z\),
\(x\leq y\) means \(\exists v\;(x+v=y)\), and one is the least positive
number.

There are effective translations in both directions.  To translate
\(W.2.A.\) to first-order arithmetic, represent an individual \(t\) by
\(p_t=2^t\) and a finite predicate \(I\) by \(y_I=F(I)\).  The atomic
translations are
\[
 \begin{array}{c|c}
 t=0&p_t=1\\
 s=t&p_s=p_t\\
 s=t+1&p_s=p_t+p_t\\
 t\in I&p_t\in_2y_I.
 \end{array}
\]
Successor terms are reduced to these clauses using fresh variables.
Translate \(\exists t\) to \(\exists p_t\,(\operatorname{Pow}(p_t)\land\cdots)\)
and \(\forall t\) to
\(\forall p_t\,(\operatorname{Pow}(p_t)\Rightarrow\cdots)\).
Finite-predicate quantifiers become unrestricted number quantifiers over
\(y_I\), because \(F\) is bijective.  Boolean connectives are unchanged.
Induction on formulas proves truth preservation under exactly these encodings.

Conversely, represent every arithmetic variable \(a\) by a finite predicate
\(I_a\) with \(a=F(I_a)\).  Equality becomes equality of predicates, and
\(a\in_2b\) becomes
\[
 \exists t\;
 \bigl(\forall s\,(s\in I_a\Longleftrightarrow s=t)\land t\in I_b\bigr).
\]
For addition \(a+b=c\), quantify a finite carry predicate \(K\).  Write
\(A_t,B_t,C_t,K_t\) for the truth values of membership at \(t\) in
\(I_a,I_b,I_c,K\), respectively, and require
\[
 \neg K_0\ \land\ \forall t\;
 \left[
  C_t\longleftrightarrow(A_t\mathbin{\oplus}B_t\mathbin{\oplus}K_t)
 \right]
 \ \land\ \forall t\;
 \left[
  K_{t+1}\longleftrightarrow
      ((A_t\land B_t)\lor(A_t\land K_t)\lor(B_t\land K_t))
 \right].
\]
Here \(\oplus\) is the Boolean exclusive-or abbreviation.  Induction on
\(t\) gives precisely binary addition with carry.  After both inputs end,
any last carry is consumed in one position, so the required carry set is
finite.  Conversely these equations force every output bit of the sum.
Eliminating compound arithmetic terms by fresh variables, and replacing
number quantifiers by finite-predicate quantifiers, completes the reverse
translation.  Thus first-order \(\mathcal B_2\)-definability and
\(W.2.A.\)-definability correspond under \(F\), and the arithmetic theory
is decidable \cite[pp.~202--203]{BruyereEtAl1994}.
Replacing every unary power-of-two test by \(\operatorname{Pow}(x)\)
also gives a decision procedure for the reduct with addition and that
unary predicate.

For \(r\geq1\), encode a tuple \((a_1,\ldots,a_r)\) by its simultaneous
binary columns, read least significant bit first.  Use the empty word for
the zero tuple; otherwise stop at the last nonzero column.  These canonical
words are exactly
\((\{0,1\}^r)^*(\{0,1\}^r\setminus\{0^r\})\cup\{\Lambda\}\).
The translations just proved and the finite-word correspondence show that a
relation on \(\mathbb N^r\) is first-order definable in \(\mathcal B_2\)
exactly when its canonical encoding language is regular.  Appending the
earlier fixed terminal marker, or removing it by right quotient, preserves
regularity.  On predicate encodings, add that marker at one past the largest
member of the union of the coordinate predicates, or at zero when all are
empty; remove it by deleting their common last point.  These operations are
definable in \(W.2.A.\) by quantifying the boundary and requiring the
corresponding membership equivalences.  Order and the largest-point condition
are definable using finite initial segments.  Thus the earlier correspondence
applies in both directions to these canonical encodings.  The recursive state
update represents feedback, and finite-set quantification supplies the
operation corresponding to Kleene star after union and concatenation
\cite[pp.~35--36, 41--42]{Buchi1959WeakReport}.

\paragraph{Elgot's controlling text and source automaton.}
Elgot's complete 1961 journal article, rather than a later compound theorem
name, controls this presentation.  The complete June 1959 Michigan report
UMR1342 supplies a primary precursor with the same section and theorem
architecture; the published 1962 errata controls every corrected formula used
below \cite{Elgot1959Report,Elgot1961,Elgot1962Errata}.

For a finite nonempty input alphabet \(I\), an \(I\)-automaton is
\[
 \mathfrak A=(S,f,d,D),
 \qquad
 f:I\times S\longrightarrow S,
 \qquad d\in S,
 \qquad D\subseteq S,
\]
where \(S\), not \(I\), is the finite nonempty set of internal states.  If
\(u=i_0\cdots i_{n-1}\), its run is the unique sequence
\[
 s_0=d,
 \qquad
 s_{k+1}=f(i_k,s_k)\quad(0\leq k<n).
\]
The behavior \(T(\mathfrak A)\) consists of the words with \(s_n\in D\); the
empty word belongs to it exactly when \(d\in D\).  Thus \(I\), \(S\), the
transition map, an input word, its run, and the accepted language have distinct
standings in the source \cite[pp.~21--22]{Elgot1961}.

\paragraph{Finite-set arithmetic as a regular language.}
Let \(x^{(0)}=x\) and \(x^{(n+1)}=(x^{(n)})'\).  Elgot's language \(L_1\) is
built from atoms \(x^{(n)}\in F_j\), \(x=y\), and \(x<y\), propositional
connectives, and individual and monadic-set quantifiers.  In the interpretation
\(L_1^1\), individuals range over \(\mathbb N\) and set variables range over
finite subsets of \(\mathbb N\).

For \(r\) finite sets \(F_1,\ldots,F_r\), put
\[
 U_r=\{0,1\}^r,
 \qquad
 c_F(n)=
 \bigl(\mathbf 1_{n\in F_1},\ldots,
       \mathbf 1_{n\in F_r}\bigr)\in U_r.
\]
For an individual term \(z\) and
\(a=(a_1,\ldots,a_r)\in U_r\), a later formula display of
\(c_F(z)=a\) is only the eliminable metalinguistic abbreviation
\[
 \bigwedge_{j:a_j=1} z\in F_j
 \quad\land\quad
 \bigwedge_{j:a_j=0} z\notin F_j.
\]
It adds no function symbol to \(L_1\).
If \(F_1\cup\cdots\cup F_r=\varnothing\), its encoding is the empty word
\(\Lambda\).  Otherwise, for
\(m=\max(F_1\cup\cdots\cup F_r)\), its encoding is
\(c_F(0)\cdots c_F(m)\).  Writing
\(0_r=(0,\ldots,0)\) and \(U_r^0=U_r\setminus\{0_r\}\), the source calls
exactly the words
\[
 U_r^*U_r^0\ \cup\ \{\Lambda\}
\]
\emph{admissible}.  For a formula \(A[F_1,\ldots,F_r]\) without free
individual variables, \(T_r(A)\) is the set of admissible encodings of its
satisfying finite-set assignments.  The formula, an assignment, its
characteristic columns, and the language \(T_r(A)\) are therefore not
interchangeable \cite[pp.~26--27]{Elgot1961}.

Theorem 5.3 gives both effective directions:
\begin{enumerate}
\item for every \(L_1^1\)-formula \(A[F_1,\ldots,F_r]\), the language
  \(T_r(A)\) is regular and a regular expression for it is effectively
  obtainable;
\item for every regular language \(\alpha\) of admissible \(U_r\)-words, there
  is a formula
  \[
   \exists G_1\cdots\exists G_s\,B[F_1,\ldots,F_r,G_1,\ldots,G_s]
  \]
  defining exactly \(\alpha\), where \(B\) has no set quantifiers and its only
  individual terms are \(x\) and \(x'\).
\end{enumerate}
For the first direction, Lemma 5.4 moves every set quantifier to the left of
every individual quantifier.  Lemma 5.5 reduces the remaining matrix to Boolean
combinations of local successor patterns.  The admissible universe is regular;
conjunction, disjunction, and negation become intersection, union, and relative
complement; and \(\exists F_r\) deletes one Boolean coordinate and then
right-truncates the resulting terminal \(0_{r-1}\)-columns.  Lemma 4.2 proves
that this truncation preserves regularity effectively.  A local formula such as
\[
 \exists x\,[c_F(x)=a\land c_F(x')=b]
\]
has an explicit regular expression, with separate terminal-zero cases.  These
steps prove closure under every formula constructor rather than merely asserting
an automata--logic analogy \cite[pp.~25--30]{Elgot1961}.

For the reverse direction, take an automaton \((S,f,d,D)\) for
\(\alpha\setminus\{\Lambda\}\), encode \(S\) by additional Boolean
coordinates, and regard \((x,y)\in U_r\times S\) as a complete letter.  The
complete-word language starts with a letter
\((x,f(x,d))\), ends in \(U_r^0\times D\), and has adjacent letters satisfying
\[
 ((x_1,y_1),(x_2,y_2))\in R
 \quad\Longleftrightarrow\quad
 f(x_2,y_1)=y_2.
\]
Elgot writes formulas for the finite start set, terminal set, and local
relation \(R\), conjoins them, and existentially projects away the state
coordinates.  The empty word is handled separately.  This supplies the
promised formula and, in particular, identifies where the reverse direction is
actually proved \cite[pp.~30--31]{Elgot1961}.

Corollary 5.8 decides truth in \(L_1^1\) by reducing a closed sentence to
emptiness or nonemptiness of its regular encoding.  Section 5.9 maps a finite
set \(F\) to
\[
 \tau(F)=\sum_{k\in F}2^k
\]
and defines addition by a finite carry set, yielding a new proof of Presburger
decidability and an effective translation.  Sections 5.10--5.16 then obtain
strictly richer decidable first-order presentations, including addition with a
power-of-two predicate, and variants for finite or cofinite sets over
\(\mathbb N\) and \(\mathbb Z\).  These are consequences of the encoding
theorem, not alternative definitions of the source automaton
\cite[pp.~31--37]{Elgot1961}.

\paragraph{Automaton design and finite-behavior synthesis.}
For finite nonempty alphabets \(A,B\), Elgot's \(A/B\)-automaton is
\[
 \mathfrak T=(S,f,d,g),
 \qquad
 f:A\times S\longrightarrow S,
 \qquad
 d\in S,
 \qquad
 g:S\longrightarrow B.
\]
Its finite behavior \(J(\mathfrak T)\subseteq (A\times B)^*\) pairs each input
word with the unique synchronous output word obtained by applying \(f\) and
then \(g\) at every position; it includes the empty paired word.  Theorem 7.1
characterizes these behaviors exactly.  A language
\(\alpha\subseteq(A\times B)^*\) is \(J(\mathfrak T)\) for some
\(\mathfrak T\) if and only if
\begin{enumerate}
\item \(\alpha\) is regular;
\item \(\alpha\) is open in the source's sense, namely closed under taking
  finite prefixes;
\item for every input word, at most one output word is paired with it; and
\item every input word has an output word paired with it.
\end{enumerate}
The proof is constructive.  It recognizes each Boolean coordinate of the
output language, takes the product of those automata, and uses functionality
and totality to recover the unique output letter
\cite[pp.~40--43]{Elgot1961}.

The operation used for design is stated before it is applied.  A paired word
\((u,v)\in\gamma\subseteq(A\times B)^*\) is \(A\)-extendable in \(\gamma\)
when, for every \(a\in A\), some \(b\in B\) makes
\((ua,vb)\in\gamma\).  It is strongly \(A\)-extendable when every prefix lies
in \(\gamma\) and is \(A\)-extendable there.  To make the effective part of
Lemma~6.1 explicit, fix a finite \((A\times B)\)-automaton
\(\mathfrak C=(S,\delta,d,D)\) recognizing \(\gamma\), so
\(\delta:(A\times B)\times S\to S\), \(d\in S\), and \(D\subseteq S\).  Put
\[
 M_0=A\times B\times D
\]
and, for \(h\geq0\), retain exactly
\[
 \begin{split}
 M_{h+1}=\{(a_0,b_0,q)\in M_h:\;&
   \forall a\in A\;\exists b\in B\;\exists q'\in D\\
 &((a,b,q')\in M_h\ \land\
   \delta((a,b),q)=q')\}.
 \end{split}
\]
This descending sequence stabilizes because \(M_0\) is finite; call its stable
value \(M_\infty\).  For each first input letter \(a\in A\), put
\[
 I_\infty(a)
 =\{b\in B:(a,b,\delta((a,b),d))\in M_\infty\}.
\]
If \(d\notin D\), or if \(I_\infty(a)=\varnothing\) for some \(a\in A\), set
\(\gamma_A=\varnothing\).  Otherwise \(\gamma_A\) contains the empty word and
the projections to \((A\times B)\)-coordinates of exactly the nonempty triple
sequences
\[
 (a_0,b_0,q_0)\cdots(a_\ell,b_\ell,q_\ell)\in M_\infty^+
\]
for which
\[
 q_0=\delta((a_0,b_0),d),
 \qquad
 q_{i+1}=\delta((a_{i+1},b_{i+1}),q_i)\quad(i<\ell).
\]
The two initial tests say exactly that the empty prefix belongs to \(\gamma\)
and is \(A\)-extendable; the stable recurrence supplies the same property at
every nonempty prefix.  This finite-state construction is therefore regular and
strongly \(A\)-extendable.  Conversely, the run triples of every strongly
\(A\)-extendable sublanguage pass the initial tests and survive every stage by
induction.  Hence the constructed language is the largest such sublanguage
\(\gamma_A\subseteq\gamma\).  This proves both termination and maximality
without introducing a separate kind of state.  Theorem 8.1 then says
\[
 \exists\mathfrak T\;J(\mathfrak T)\subseteq\gamma
 \quad\Longleftrightarrow\quad
 \gamma_A\ne\varnothing.
\]
Fix a total order on the finite alphabet \(B\), and give equal-length
\(B\)-words its induced lexicographic order.  When \(\gamma_A\) is nonempty,
retain within \(\gamma_A\) the least compatible output for each input.
Lemma~6.5 preserves regularity; strong \(A\)-extendability supplies prefix
closure and input-totality; and the selection is functional.  Theorem~7.1
therefore produces the required automaton.  Corollary 8.2 decides whether all automata
whose behavior lies in \(\gamma\) have the same behavior: the condition is
vacuously true when \(\gamma_A=\varnothing\), and otherwise it compares the
selections obtained from the fixed order on \(B\) and its reverse.  Section 8.3 encodes
an \(A/B\)-automaton itself by an \(L_1^1\)-formula, so behavior containment is
also expressible as a truth question \cite[pp.~37--46]{Elgot1961}.

For infinite paired sequences, if \(\alpha^L\) denotes all finite prefixes of
members of \(\alpha\subseteq(A\times B)^{\mathbb N}\), Theorem 7.3 says that
\(\alpha\) is the infinite behavior of an \(A/B\)-automaton exactly when
\(\alpha^L\) is regular, equal input prefixes force equal output prefixes, and
every infinite input has an output.  This theorem is about causal synchronous
behavior; it is not an acceptance condition for B\"uchi automata
\cite[pp.~43--44]{Elgot1961}.

\paragraph{The infinite synthesis reduction.}
Section 8.6 changes the interpretation of set variables from finite to
arbitrary subsets of \(\mathbb N\).  Its restricted specification has the form
\[
\exists\mathbf s\,\forall x\,
 M(\mathbf a,\mathbf b,\mathbf s,x),
\]
where
\[
 \mathbf a=(a_1,\ldots,a_m),\qquad
 \mathbf b=(b_1,\ldots,b_{m'}),\qquad
 \mathbf s=(s_1,\ldots,s_r)
\]
are tuples of monadic predicate variables, \(m,m',r\geq0\), and
\(d=m+m'+r\).  For any formula \(C\) with those \(d\) predicate variables
and any assignment \(\mathbf P=(P_1,\ldots,P_d)\in\mathcal P(\mathbb N)^d\),
put
\[
 c_{\mathbf P}(t)
   =\bigl(\mathbf 1_{t\in P_1},\ldots,
           \mathbf 1_{t\in P_d}\bigr)
   \in\{0,1\}^d
\]
and define
\[
 T_d^\infty(C)
 =\bigl\{c_{\mathbf P}\in(\{0,1\}^d)^{\mathbb N}:
          \mathbf P\text{ satisfies }C\bigr\}.
\]
Thus \(T_d^\infty(C)\) is a language of characteristic-column streams, not
the formula, an assignment, or an unencoded model.  The matrix \(M\) is
quantifier-free, has only the displayed variables free, may
contain numerals, and contains neither equality nor order.  If
\(z^{(h)}\) denotes the \(h\)-fold successor of \(z\), the individual terms
relevant to the construction are \(x^{(h)}\) and \(0^{(h)}\); put
\[
 n=1+\max\bigl(\{0\}\cup
       \{h:x^{(h)}\text{ or }0^{(h)}\text{ occurs in }M\}\bigr).
\]
These restrictions make truth depend on a bounded window of characteristic
columns and the fixed initial columns named by numerals.  The unprojected
model language is a finite union.  Writing \(\bigwedge_x M\) for the
universal quantifier,
\[
 T^\infty_d(\bigwedge_x M)
 =S_1\cup\cdots\cup S_k,
\]
where each \(R_j\) is a finite \(n\)-ary relation.  Its field is the union of
its coordinate projections,
\[
 A_j=\operatorname{field}(R_j)
 =\{a:\exists(a_0,\ldots,a_{n-1})\in R_j\;
             \exists h<n\;(a=a_h)\}.
\]
An infinite \(R_j\)-sequence is an
\(f\in A_j^{\mathbb N}\) satisfying
\((f(t),\ldots,f(t+n-1))\in R_j\) for every \(t\), and
\(S_j\) consists exactly of those sequences whose length-\(n\) prefix lies in
a finite set \(E_j\subseteq A_j^n\).
Construct the union by considering each possible initial word
\(b\in U_d^n\).  Interpret terms \(0^{(h)}\) using \(b_h\), and
terms \(x^{(h)}\) using coordinate \(h\) of a moving length-\(n\) window.
Let \(R_b\subseteq U_d^n\) contain exactly the windows satisfying this
propositional matrix, and set \(E_b=\{b\}\cap R_b\).
A model has one initial word \(b\), and all its windows belong to
\(R_b\).  Conversely, an admitted stream has that initial word and
satisfies \(M\) at every position.  Taking all such \(b\) proves the
finite union.  If every component is empty, the union is the empty language.

Reduce the finite union to a projection of one relation language as follows.
For \(n\geq2\), replace the field of each
\(R_j\) by its disjoint tagged copy
\(\operatorname{field}(R_j)\times\{j\}\) and define the \(n\)-ary relation
\[
 R'_j
 =
 \left\{
   ((a_0,j),\ldots,(a_{n-1},j)):
   (a_0,\ldots,a_{n-1})\in R_j
 \right\}.
\]
Tag every coordinate of every word in \(E_j\) in the same way, then take the
unions of the tagged relations and prefix sets.  Consecutive length-\(n\)
windows overlap, so every admitted infinite tagged sequence has one constant
tag \(j\); projection \((a,j)\mapsto a\) therefore yields exactly
\(S_1\cup\cdots\cup S_k\).

If \(n=1\), then \(R_j\subseteq A_j\) and \(E_j\subseteq A_j\).  Replace
the unary constraint by
\[
 \widehat R_j=R_j\times R_j,
 \qquad
 \widehat E_j=(E_j\cap R_j)\times R_j\subseteq A_j^2.
\]
The binary windows require every letter to lie in \(R_j\), while
\(\widehat E_j\) retains exactly the original first-letter condition.  Thus
this defines the same infinite \(R_j\)-sequences, after which the tagged binary
construction forces a constant tag.  Both arity cases therefore reduce to
one finite relation language followed by projection.  Section~8.6.3 obtains
an effective infinite-behavior synthesis procedure by forming the regular
language of finite prefixes and applying Theorem 8.1
\cite[pp.~44--47]{Elgot1961}
\cite[pp.~558--559]{Elgot1962Errata}.

\paragraph{Elgot's negative boundary.}
Chapter IV adds addition to successor and first-order quantification over a
language with monadic predicates.  Lemma 10.1 encodes proofs in Post normal
systems.  Theorem 10.2 concludes that, when the predicates range over finite
subsets of \(\mathbb N\), the satisfiable formulas are recursively enumerable
but undecidable, with maximal recursively enumerable degree.  Sections
10.3--10.5 transfer the negative result to unrestricted predicates, integers,
and ultimately periodic predicates.  Corollary 10.6 makes the design boundary
especially sharp: for a fixed input-free automaton, satisfaction of a given
formula is decidable, while existence of some automaton satisfying the formula
is undecidable.  Elgot's positive algorithms therefore depend on the exact
language and specification form established earlier; they are not a blanket
decidability claim for arithmetic automaton design
\cite[pp.~47--50]{Elgot1961}.

\paragraph{Trakhtenbrot's two source bodies and formal system.}
The registered source is Trakhtenbrot's four-page 1961 announcement.  His
twenty-nine-page 1962 article gives the complete definitions, proofs, examples,
and qualifications compressed there.  They have the same Russian title but
different official Math-Net English renderings:
\enquote{Single-Place Predicates} in 1961 and
\enquote{One-Place Predicates} in 1962.  The announcement controls claims
made under its theorem numbers and pp.~326--329; the expansion controls the
five-section development on pp.~103--131.  They are related primary texts, not
silently interchangeable versions
\cite{Trakhtenbrot1961,Trakhtenbrot1962Expanded}.

The source prints the name of its interpreted formal system as the Cyrillic
capital letter И; write it here as \(\mathcal I\).  Individual variables range
over
\[
  \mathbb N_{+}=\{1,2,3,\ldots\},
\]
monadic predicate variables range over all subsets of \(\mathbb N_{+}\), and
immediate successor is primitive.  For an individual variable \(t\), put
\[
 t^{(0)}=t,
 \qquad
 t^{(k+1)}=(t^{(k)})'
 \quad(k\geq0).
\]
The atomic formulae are \(X(t^{(k)})\), where \(X\) is a monadic predicate
variable; Boolean operations, individual quantifiers, and monadic-predicate
quantifiers generate all formulae.  Equality, order, bounded individual
quantifiers, constants, and terms \(t+k\) are explicitly eliminable
conveniences, not additional primitives.  Thus \(\mathcal I\) names the formal
language with this grammar and interpretation; it is neither an automaton nor
the weak finite-set system used by B\"uchi.

A free predicate is \emph{fictitious outside} an interval when changing its
values outside that interval cannot change the formula's truth.  This is
semantic independence, not a test for whether a symbol occurs in the written
formula.  Fix a formula whose only free individual variable is \(t\).  It is
\emph{\(t\)-controlled from above} when every individual quantifier over
\(\sigma\) has a range of the form
\(\sigma<\mu\) or \(\lambda<\sigma<\mu\), and it is
\emph{\(t\)-controlled from below} when every such range has the form
\(\lambda<\sigma\) or \(\lambda<\sigma<\mu\), where \(\lambda,\mu\) are
terms.  An \emph{upper \(t\)-formula} is a \(t\)-controlled-from-below formula
in which every atomic argument is an individual variable or a variable plus a
numeral constant, never a bare constant, and every individual-quantifier range
has the form \(t<\sigma\) or \(t<\sigma<\mu\).  A
\emph{lower \(t\)-formula} is dually \(t\)-controlled from above, has the same
atomic-argument restriction, and has ranges
\(\sigma<t\) or \(\lambda<\sigma<t\).  These are syntactic classes.  As their
separate semantic consequences, an upper \(t\)-formula is fictitious outside
\([t,\infty)\), while a lower \(t\)-formula is fictitious outside
\([1,t+s]\), where \(s\) is the largest numeral offset occurring in an
argument of a free predicate.  Theorem~1 of the announcement, expanded in
\S1.4, effectively turns every \(\mathcal I\)-formula into a Boolean
combination of lower and upper \(t\)-formulae.  Its proof moves predicate
quantifiers outside individual quantifiers, pushes the individual quantifiers
inward, splits them at \(t\), and distributes each remaining predicate
quantifier across the independent prefix and tail parts.  The finitely many
lower-window possibilities are the source of the finite constructions below
\cite[pp.~326--327]{Trakhtenbrot1961}
\cite[pp.~103--109]{Trakhtenbrot1962Expanded}.

\paragraph{Anticipation, memory, and finite operators.}
For finite nonempty alphabets \(A,B\), an operator is a total map
\[
  T:A^{\mathbb N_{+}}\longrightarrow B^{\mathbb N_{+}},
  \qquad a\longmapsto b.
\]
It has anticipation \(p\geq0\) when agreement of two inputs through position
\(t+p\) forces their outputs to agree at position \(t\).  Anticipation zero is
therefore exactly the prefix-preservation condition: extending an input word
cannot change the output word already produced.

Finite memory is a different restriction.  It means that there are a finite
state set \(Q\), an initial state \(q(1)\), and maps
\[
 \Phi:Q\times A^{\mathbb N_{+}}\to B,
 \qquad
 \Psi:Q\times A^{\mathbb N_{+}}\to Q
\]
such that, with \(a_{\geq t}=(a(t),a(t+1),\ldots)\),
\[
  b(t)=\Phi(q(t),a_{\geq t}),
  \qquad
  q(t+1)=\Psi(q(t),a_{\geq t}).
\]
The least possible \(\lvert Q\rvert\) is the operator's \emph{weight}.  Future
dependence may still be unbounded in this definition: Trakhtenbrot's
one-state example asks whether some future input letter is \(1\), and need not
be recursive.

A \emph{finite operator} has both finite memory and bounded anticipation.  For
some fixed \(p\), it can be written in the canonical finite-window form
\[
  b(t)=\Phi\bigl(q(t),a(t),\ldots,a(t+p)\bigr),
  \qquad
  q(t+1)=\Psi\bigl(q(t),a(t),\ldots,a(t+p)\bigr).
\]
The 1961 announcement calls the finite operators with \(p=0\)
\enquote{Л-operators}; the 1962 expansion calls the same automaton-realizable
class \(A\)-operators.  Whenever a sentence below combines the two versions,
this presentation expressly adopts the expansion's Latin \(A\) as normalized
terminology; it does not attribute that letter to the announcement.
In the hyphenated name, the letter labels the automaton-realizable class; it is
not a free occurrence of the input-alphabet variable.  Zero anticipation
alone, finite memory alone, finite operator, and \(A\)-operator are consequently
four different conditions.

The expansion calls a stream \(f\) \enquote{periodic} when some finite preperiod
\(\alpha\) and period \(\beta>0\) satisfy
\(f(\alpha+t)=f(\alpha+\beta+t)\) thereafter; in current terminology this is
\emph{ultimately periodic}.  Section~2 proves that primitive-recursive schemes
are too broad as a synthesis language: deciding whether such a scheme defines
an ultimately periodic output, or extracting canonical finite-state equations
when it does, is undecidable
\cite[pp.~109--113]{Trakhtenbrot1962Expanded}.

\paragraph{Portions, homogeneity, and definability.}
Identify a Boolean predicate \(X\) with the Cantor-space point
\(X(1)X(2)\cdots\).  If \(M\) is a set of such points and \(u\) is a finite
Boolean word, define its portion after \(u\) by
\[
  M_u=\{Z:uZ\in M\}.
\]
The 1961 announcement explicitly permits this set to be empty.  The 1962
expansion instead introduces a portion only after requiring the corresponding
prefix intersection to be nonempty.  This handbook uses the announcement's
all-prefix convention.  For finite alphabets \(A_1,\ldots,A_n\),
\(M\subseteq\prod_{i=1}^n A_i^{\mathbb N_{+}}\), an integer \(r\geq0\), and
equal-length coordinate prefixes \(u_i\in A_i^r\), the tuple portion is exactly
\[
 M_{(u_1,\ldots,u_n)}
 =\bigl\{(Z_1,\ldots,Z_n)\in\prod_{i=1}^n A_i^{\mathbb N_{+}}:
        (u_1Z_1,\ldots,u_nZ_n)\in M\bigr\}.
\]
Thus concatenation and shifting are coordinatewise, and the coordinates are
not interleaved.  The set \(M\) is \emph{homogeneous of degree at most \(k\)}
when the family \(\{M_u:u\text{ finite}\}\) has at most \(k\) distinct
members.  A finite list of all distinct portions is a basis; appending a
letter carries one basis member to another and so determines its internal
transition map.  These are properties and finite data of a set of streams,
not new automaton species.

Here and below \emph{closed} means closed in the Cantor product topology.
Let \(A,B\) be finite nonempty alphabets.  The graph of a total operator
\(T:A^{\mathbb N_{+}}\to B^{\mathbb N_{+}}\) is
\[
 G_T=\{(X,T(X)):X\in A^{\mathbb N_{+}}\}.
\]
If \(T\) has zero anticipation, \(u\in A^r\), and \(Z\in A^{\mathbb N_{+}}\),
then the length-\(r\) prefix of \(T(uZ)\) is independent of \(Z\); call it
\(v_u\in B^r\).  Define
\[
 \operatorname{tail}_r(Y)(t)=Y(t+r),
 \qquad
 T_u:A^{\mathbb N_{+}}\longrightarrow B^{\mathbb N_{+}},
 \qquad
 T_u(Z)=\operatorname{tail}_r(T(uZ)).
\]
For every \(v\in B^r\), the typed tuple definition gives
\[
 (G_T)_{(u,v)}
 =\begin{cases}
   G_{T_u},&v=v_u,\\
   \varnothing,&v\neq v_u.
  \end{cases}
\]
Define the family of residual maps and the empty-portion indicator by
\[
 \mathcal R_T=\{T_u:u\in A^*\},\qquad
 \epsilon_T=
 \begin{cases}
  1,&\exists r\geq0\;\exists u\in A^r\;\exists v\in B^r\;(v\ne v_u),\\
  0,&\text{otherwise}.
 \end{cases}
\]
Then a zero-anticipation graph \(G_T\) is closed, and it is homogeneous
if and only if \(\mathcal R_T\) is finite.  When finite, its degree in
the all-prefix convention is exactly
\[
 \deg(G_T)=|\mathcal R_T|+\epsilon_T.
\]
For closedness, take \((X,Y)\notin G_T\) and choose a position where
\(Y\) differs from \(T(X)\).  Fix both prefixes through that position.
Zero anticipation fixes the corresponding output prefix, so the resulting
product cylinder is disjoint from \(G_T\).  Every point outside the graph
therefore has an open neighborhood outside it, proving closedness.
For the degree statement, the displayed portion identity shows that every
portion is some \(G_{T_u}\) or the empty set.  Conversely, every
\(G_{T_u}\) occurs after the consistent prefix \((u,v_u)\).
All these graphs are nonempty because their maps are total on the nonempty
space \(A^{\mathbb N_+}\), and two such graphs agree exactly when their
maps agree.  The empty set is a further distinct portion exactly when
\(\epsilon_T=1\).  This proves both the equivalence and the exact degree.

Memory weight with whole-future updates need not bound this residual degree.
For example, take \(A=B=\{0,1\}\) and consider
\[
 T(X)(t)=
 \begin{cases}
  X(t),&X(1)=1,\\
  0,&X(1)=0.
 \end{cases}
\]
Its output through time \(t\) depends only on the input through time \(t\), so
its anticipation is zero.  It has weight exactly two.  A two-state
finite-memory presentation is obtained from
\(Q=\{0,1\}\), \(q(1)=0\), and, for
\(z=(z_0,z_1,\ldots)\),
\[
 \Phi(q,z)=(1-q)z_0,
 \qquad
 \Psi(q,z)=q\mathbin{\mathsf{xor}}z_0\mathbin{\mathsf{xor}}z_1.
\]
Indeed, \(q(t)=X(1)\mathbin{\mathsf{xor}}X(t)\).  A one-state presentation
would use the same suffix function at every time; its value at time one would
have to be \(z_0\), contradicting the required time-two output on an input
beginning \(01\).  Yet \(G_T\), its portions after the paired one-letter
coordinate prefixes \((0,0)\) and \((1,1)\), and its portion after
\((0,1)\) are respectively the original graph, the constant-zero graph, the
identity graph, and \(\varnothing\).  Every nonempty input prefix fixes
the first input bit, so the remaining map is either constant zero or the
identity.  Thus there are exactly three residual maps and one empty portion,
giving degree four.  The memory state in the displayed presentation depends
on both the first and current input bits; its transition function may inspect
the future suffix.  Counting those memory states does not count residual maps.
A singleton is homogeneous exactly when its unique stream is ultimately
periodic.  Homogeneity alone does not imply \(\mathcal I\)-definability: there
are uncountably many homogeneous sets of degree at most \(1\), but only
countably many formulae.

The announcement's Theorems~2--3 and expanded \S3 give the exact positive
claims.  Every \(\mathcal I\)-definable set is homogeneous, and a defining
formula effectively determines the finite internal transition map among its
portions.  Conversely, every \emph{closed} homogeneous set is
\(\mathcal I\)-definable, by a formula of the recurrence shape
\[
 \begin{aligned}
  \varphi(X)\ \equiv
  (\exists q_1\cdots\exists q_k\exists U)\bigl[&
    q_1(1)=1\ \land\ q_2(1)=\cdots=q_k(1)=0\\
    &{}\land(\forall t)\,
    \theta\bigl(X(t),q_1(t),\ldots,q_k(t),U(t),
                 q_1(t+1),\ldots,q_k(t+1)\bigr)
  \bigr].
 \end{aligned}
\]
Here \(X\) is the free predicate whose membership in the set is defined;
\(q_1,\ldots,q_k\) are bound monadic predicates recording the finite basis
state; \(U\) is a bound monadic choice predicate; and \(\theta\) is
quantifier-free.  The literal initial values select the basis member
corresponding to the original set, rather than an unspecified union of
portions.  For a nonempty homogeneous set \(M\), list its nonempty basis
portions as \(P_1=M,P_2,\ldots,P_k\).  Appending a bit \(b\in\{0,1\}\) at
\(P_i\) is permitted exactly when the successor portion \((P_i)_b\) is
nonempty; that successor is then one of the listed basis members.  Given a
choice stream \(U\in\{0,1\}^{\mathbb N_+}\), start at \(P_1\), output the
only permitted bit when there is one and output \(U(t)\) when both bits are
permitted, and move to the corresponding successor portion.  Every nonempty
portion permits at least one bit, so this defines a total map
\(\{0,1\}^{\mathbb N_+}\to\{0,1\}^{\mathbb N_+}\); its image is exactly the
closure \(\overline M\), because its outputs are precisely the streams whose every
finite prefix extends to a member of \(M\).  If \(M\) is empty, no total map
can have that image.  The empty set is instead handled in the same recurrent
syntactic shape by conjoining the displayed initial clause \(q_1(1)=1\) with
\(q_1(1)=0\).  This proves the empty case as well.  Every formula of
the displayed recurrent shape also defines a closed homogeneous set.
Consequently the \(\mathcal I\)-definable sets form the
least class containing the closed homogeneous sets and closed under
complement, finite union, and predicate projection.  The announcement leaves
the exact represented Borel and projective classes open
\cite[pp.~326--328]{Trakhtenbrot1961}
\cite[pp.~113--119]{Trakhtenbrot1962Expanded}.

There is an important effectiveness qualification.  The internal transition
map follows directly from the finite family of formula portions.  Choosing an
output requires deciding which portions are nonempty.  Trakhtenbrot therefore
calls a procedure which assumes a truth algorithm for closed
\(\mathcal I\)-sentences a \emph{conditional \(\mathcal I\)-algorithm}.  The
condition is not erased by later knowledge.  The source notes that Putnam's
decision result makes the construction unconditional for the displayed
existential-predicate recurrence form.

\paragraph{Finite realization.}
Suppose an \(\mathcal I\)-formula \(\varphi(X,Y)\) defines the graph of a total
operator.  Theorem~4 of the announcement and expanded \S4, Theorem~1, say that
if this operator is continuous, then it is finite and its weight is bounded by
the homogeneity degree of the graph.  Compactness makes continuity uniform on
each of the finitely many nonempty portions, giving one anticipation bound;
the portion basis supplies finite memory.  If the operator has zero
anticipation, it belongs to the class called a Л-operator in the announcement
and an \(A\)-operator in the expansion.  Expanded \S4, Theorem~2, deliberately
drops continuity and obtains only finite memory, not a finite operator; no
bounded anticipation follows in that case.

For any relation specified by \(\varphi(X,Y)\), an operator \(T\) is a
\emph{solution} when \(\varphi(X,T(X))\) holds for every input \(X\), equivalently
when the graph of \(T\) is contained in that relation.  This is weaker than
saying that \(\varphi\) defines exactly the graph of \(T\), and the distinction
is needed immediately.  The existence form, Theorem~\(4'\), concerns a closed
homogeneous relation rather than an already defined graph.  If it has a
continuous operator solution, it has a finite solution of bounded weight; if
it has any zero-anticipation solution, it has a Л-operator solution---an
\(A\)-operator in the expansion's terminology---with the same bound.  The
announcement separately records the analogous assertion for an open relation,
but without any bound on the weight of the resulting solution.

Trakhtenbrot strengthens existence to a precise parametrization.  In the
expansion's notation, an \(A\)-operator \(Y=T(X,U)\), where \(U\) is a finite
tuple of parameter streams, is a \emph{general \(A\)-solution} of
\(\varphi(X,Y)\) when both of the following hold; the announcement prints the
same definition with Л in place of \(A\):
\begin{enumerate}
\item every \(A\)-operator solution \(Y=T_0(X)\) is obtained as
  \(T_0(X)=T(X,S(X))\) for some \(A\)-operator tuple \(S\); and
\item every \(A\)-operator substitution \(U=S(X)\) yields a solution.
\end{enumerate}
Replacing \(A\)-operators everywhere by arbitrary zero-anticipation operators
defines the source's general zero-anticipation solution.  The adjective
\enquote{general} therefore means
\enquote{parametrizes all and only the solutions of the named class}, not
merely \enquote{is one solution}.

Theorem~5 of the announcement and expanded \S4, Theorem~3, state that for the
closed relation defined by an \(\mathcal I\)-formula---and therefore homogeneous
of some finite degree---either no zero-anticipation solution exists, or a
general solution of the automaton-realizable class exists with weight at most
that degree and is simultaneously a general zero-anticipation solution.  This
is a general Л-solution in the announcement and a general \(A\)-solution in the
expansion.  The proof repeatedly deletes empty or nonextendible portions.  If
the initial portion survives, free parameter bits select among the remaining
outputs and the finite portion transition supplies memory; closedness turns
compatibility of every finite prefix into membership in the relation.  The
arbitrary closed-\(\mathcal I\)-formula construction is conditional because
deletion asks whether a portion is empty.  Expanded Theorem~\(3'\) gives the
exact fixed-lookahead extension for the same closed \(\mathcal I\)-definable
relation: for each fixed \(p\), either no solution whose anticipation is
bounded by \(p\) exists, or one finite operator with that bound parametrizes
all solutions with the same bound.  It is not one unbounded parametrization
for every \(p\) at once
\cite[pp.~328--329]{Trakhtenbrot1961}
\cite[pp.~119--124]{Trakhtenbrot1962Expanded}.

\paragraph{Explicit synthesis and the finite-word boundary.}
Section~5 calls
\[
  Y(t)\ \longleftrightarrow\ \varphi(X_1,\ldots,X_m;t)
\]
an explicit definition.  It always defines a finite-memory operator.  It
defines a finite operator of anticipation \(p\) exactly when its truth is
independent of input-predicate values outside \([1,t+p]\), and it defines an
\(A\)-operator exactly when the independence interval is \([1,t]\).  For the
zero-shift prefix-controlled case, an unconditional construction fixes the
first input column, records the resulting shifted formula, and repeats.  The
normal form permits only finitely many resulting formulae, so they become the
states and determine the output and transition tables.  The worked example on
pp.~126--127 has two such formulae and prints both tables.

For the source's positive word lengths, every zero-anticipation stream operator
induces typed word maps
\[
  T_t:A^t\longrightarrow B^t,
  \qquad
  T_t(u)=T(X)\mathbin{\upharpoonright}\{1,\ldots,t\}\quad(t\geq1),
\]
where \(X\) is any infinite extension of \(u\).  Two such extensions agree
through \(t\), so zero anticipation proves that \(T_t\) is independent of the
chosen extension.  It also proves prefix compatibility:
\(T_s(u\mathbin{\upharpoonright}\{1,\ldots,s\})
=T_t(u)\mathbin{\upharpoonright}\{1,\ldots,s\}\) for \(s\leq t\).

The article states the following definitions for binary words and says that
arbitrary finite nonempty alphabets are analogous.  A language
\(L\subseteq\bigcup_{t\geq1}B^t\) is \emph{automaton-enumerable} when some
\(A\)-operator \(T\)
satisfies
\[
  L\cap B^t=T_t(A^t)\qquad\text{for every }t\geq1.
\]
For an \(A\)-operator
\(R:B^{\mathbb N_{+}}\to\{0,1\}^{\mathbb N_{+}}\), the terminal result on
\(v\in B^t\) is \(R(X)(t)\), where \(X\) is any infinite extension of \(v\);
again zero anticipation makes this well-defined.  The language is
\emph{automaton-decidable} when such an \(R\) has terminal result \(1\) exactly
for the words in \(L\).  Every automaton-enumerable language is
automaton-decidable.  Conversely, an automaton-decidable language \(L\)
is automaton-enumerable under the following three hypotheses:
\[
 \begin{gathered}
  L\neq\varnothing,\\
  \forall v\in L\;\forall u\in\bigcup_{t\geq1}B^t\;
     (u\text{ a prefix of }v\ \Longrightarrow\ u\in L),\\
  \forall v\in L\;\exists w\in L\;
     (v\text{ is a strict prefix of }w).
 \end{gathered}
\]
From any member, the nonempty-prefix clause supplies a word of length one;
the strict-extension clause followed by another nonempty prefix then supplies
each successive positive length needed by the enumeration construction.
The empty language is automaton-decidable by a constant-zero operator,
but a total length-preserving operator has a nonempty range at every
positive length and therefore cannot enumerate it.

Finally the weak system \(\mathcal I'\) restricts predicate variables to
predicates that are eventually zero.  For such a predicate \(X\), the finite
word \(X(1)\cdots X(t)\) is a \emph{special beginning} exactly when
\(X(\tau)=0\) for every \(\tau>t\); it is therefore not an arbitrary finite
prefix.  Expanded \S5, Theorem~3, asserts that for every \(\mathcal I\)-formula
there exists an \(A\)-operator recognizing exactly the special beginnings of
the satisfying eventually-zero predicates.  For an \(\mathcal I'\)-formula
the operator is effectively constructible, a conclusion the source explicitly
attributes to B\"uchi
\cite[pp.~124--131]{Trakhtenbrot1962Expanded}.

This is Trakhtenbrot's literal finite-word bridge.  The familiar sentence
\enquote{finite-word MSO-definable if and only if regular} is not printed as
the four-page article's theorem.  His source package is broader: arbitrary
predicate quantification, finite memory and anticipation, Cantor-space
topology, definability, recognition, synthesis, and general solutions.  The
1961 announcement also reports its conditional-algorithm boundary and the
author's comparisons with Church, B\"uchi, and Elgot as historical claims,
not as an independently adjudicated priority theorem.

B\"uchi's report presentation is source-controlled at its 1959 precursor
strength, including the printed defects and explicitly owned repairs above;
literal reconciliation with the 1960 journal version remains separate.
Elgot's article, precursor, and published errata are source-controlled
independently.  Trakhtenbrot's two Russian bodies are reconstructed here at
their distinct announcement and expanded-article strengths; no claim of a
complete line-by-line translation is made.

At the level of the modern construction, Thomas's 1996 report precursor makes
the finite-word route exact.  Let \(Q\) be a finite nonempty state set,
\(\Sigma\) a finite alphabet, \(q_0\in Q\), and \(F\subseteq Q\).  Renaming
only the source alphabet and transition relation, Section~3.1 uses
\[
 \mathcal A=(Q,\Sigma,q_0,\Delta,F),
 \qquad
 \Delta\subseteq Q\times\Sigma\times Q.
\]
For \(w=a_0\cdots a_{n-1}\in\Sigma^*\), put
\[
 \operatorname{Pos}(w)=\{i\in\mathbb N:i<n\},
 \qquad
 S^w(i,j)\Longleftrightarrow j=i+1<n,
 \qquad
 P_a^w=\{i<n:a_i=a\}\quad(a\in\Sigma).
\]
Thomas's finite-word interpretation of \(\operatorname{MSO}[S]\) is the
structure
\[
 \mathfrak W_w=
 \bigl(\operatorname{Pos}(w),S^w,(P_a^w)_{a\in\Sigma}\bigr).
\]
Its individual variables range over \(\operatorname{Pos}(w)\), and its monadic
variables range over \(\mathcal P(\operatorname{Pos}(w))\).  A successful run
of \(\mathcal A\) on \(w\) is a sequence \(r(0),\ldots,r(n)\) with
\(r(0)=q_0\),
\((r(i),a_i,r(i+1))\in\Delta\) for \(i<n\), and \(r(n)\in F\).
Theorem~3.1 states that finite-word languages recognized by such finite automata are
exactly the \(\operatorname{MSO}[S]\)-definable languages, effectively in both
directions \cite[\S3.1]{Thomas1996Report}.

For the automaton-to-formula direction, disjoint monadic tracks encode
\(r(0),\ldots,r(n-1)\); the final input letter is treated separately by requiring
its tracked state to have an \(F\)-successor.  The report explicitly notes that
this displayed formula is true on the empty word when all tracks are empty, so
an existence guard must be conjoined when \(\mathcal A\) rejects the empty word.
For the converse, Thomas first passes to the singleton-set version
\(\operatorname{MSO}_0\), in which every individual variable is replaced by a
monadic variable constrained to be a singleton.  For
\(w=a_0\cdots a_{n-1}\), the predicates \(P_a^w\) above are the source's
letter predicates \(Q_a\), renamed to avoid collision with the state set
\(Q\).  For subsets \(X,Y\subseteq\operatorname{Pos}(w)\), the remaining
atomic predicates have the exact meanings
\[
 \operatorname{Sing}(X)\Longleftrightarrow
   \exists i<n\;(X=\{i\}),
 \qquad
 \operatorname{Suc}(X,Y)\Longleftrightarrow
   \exists i<n\;\exists j<n\;
      (X=\{i\}\land Y=\{j\}\land j=i+1).
\]
Automata are supplied for the atomic forms
\(X_j\subseteq X_k\), \(\operatorname{Sing}(X_j)\),
\(\operatorname{Suc}(X_j,X_k)\), and \(X_j\subseteq P_a^w\); complement and union
interpret negation and disjunction, while nondeterministic projection of one
bit track interprets existential set quantification.  The same reduction sends
a closed sentence of \(WS1S\), the weak monadic second-order theory of one
successor on \(\mathbb N\), to an input-free finite automaton.  Here individual
variables range over \(\mathbb N\), monadic variables range over finite subsets
of \(\mathbb N\), and successor has its usual meaning.  Thus \(WS1S\) is not
the finite structure \(\mathfrak W_w\); its truth is decided by successful-run
existence.
Reversed binary words and an auxiliary finite carry set then give the report's
Presburger route
\cite[pp.~8--12]{Thomas1996Report}.

This modern proof architecture is not substituted for B\"uchi's different
special-predicate encoding, Elgot's design calculus, or Trakhtenbrot's broader
operator-and-predicate theory.  It closes the 1996-report locator only; it does
not turn the survey into a transcription of the unreconciled 1960 B\"uchi
article or either Trakhtenbrot source, and it does not replace the separately
reconstructed and erratum-corrected 1961 Elgot presentation.

\subsubsection{Rabin's successor structures and ordinary tree automata}

Rabin's 1969 article \emph{Decidability of Second-Order Theories and Automata
on Infinite Trees} controls the historical definitions, numbered results, and
proof architecture below; the later Thomas route is kept separate at the end.

For any set \(A\), Rabin writes \(A^*\) for the finite words over \(A\), uses
juxtaposition for concatenation, and writes \(\Lambda\) for the empty word.
The source symbol \(\Lambda\) is retained below; no editorial
\(\epsilon\) is substituted for it.
To distinguish the paper's two order symbols typographically, write
\[
 x\preccurlyeq_{\mathrm{is}}y
   \quad\Longleftrightarrow\quad
   \exists z\in A^*\;(y=xz)
\]
for the initial-segment order and write
\(\preccurlyeq_{\mathrm{lex}}\) for the lexicographic order induced by the
given order on \(A\).  These subscripts are editorial disambiguations, not
additional relations.  For \(a\in A\), let
\[
 r_a:A^*\longrightarrow A^*,
 \qquad r_a(x)=xa.
\]
The predecessor used below is the map
\[
 \operatorname{pd}:A^*\longrightarrow A^*,
 \qquad
 \operatorname{pd}(\Lambda)=\Lambda,
 \qquad
 \operatorname{pd}(xa)=x
 \quad(x\in A^*,\ a\in A).
\]
Definition 1.1 applies to every \(0<n\leq\omega\):
\[
 T_n=n^*=\{i:i<n\}^*,
 \qquad
 \mathfrak N_n
  =\bigl\langle T_n,(r_i)_{i<n},
      \preccurlyeq_{\mathrm{is}},\preccurlyeq_{\mathrm{lex}}\bigr\rangle.
\]
The monadic second-order theory of this structure is \(SnS\).  When \(n\) is
finite, both orders are monadic-second-order definable from the successors;
the displayed structure retains them because that elimination fails for
\(T_\omega\).  Theorem 1.1 states that \(S2S\) is decidable.  Rabin says that
the paper's procedures are elementary recursive in Kalmar's sense
\cite[pp.~1, 4]{Rabin1969}.  That complexity claim is not retained here:
Meyer's later ICM report states a non-elementary tower lower bound and credits
the two-successors-and-prefix case to Meyer and Stockmeyer; that theory is a
sublanguage of the displayed binary theory
\cite[pp.~478--479]{Meyer1975Complexity}.  The later lower bound leaves Rabin's
decidability and effectiveness result intact but rules out the stated
elementary-recursive worst-case bound.  Theorem 1.12 gives a definable
interpretation of \(\mathfrak N_\omega\) in the binary tree and thereby also
decides \(S\omega S\)
\cite[pp.~3--5, 10--11]{Rabin1969}.

Put \(T=T_2=\{0,1\}^*\).  For \(x\in T\), the subtree rooted at \(x\) is
\[
 T_x=\{y\in T:x\preccurlyeq_{\mathrm{is}}y\}.
\]
A path through \(T_x\) contains \(x\), contains exactly one of \(y0,y1\)
whenever it contains \(y\), and contains the immediate predecessor of each of
its nodes other than \(x\), namely \(\operatorname{pd}(y)\).  Let \(\Sigma\) be a finite alphabet.  A
\(\Sigma\)-valued tree is a pair \((v,T_x)\) with \(v:T_x\to\Sigma\).
The set of all such valued subtrees as \(x\) varies is \(V_\Sigma\).
Definition 1.3 distinguishes a transition table
\[
 \mathfrak A_0=\langle S,M\rangle,
 \qquad
 M:S\times\Sigma\longrightarrow\mathcal P(S\times S),
\]
from an automaton
\[
 \mathfrak A=\langle S,M,S_0,F\rangle,
 \qquad S_0\subseteq S,
 \qquad F\subseteq\mathcal P(S).
\]
The finite set \(S\) is the state set, and the move table need not be total.
For any finite set \(R\), any set \(X\), and any map \(\phi:X\to R\), define
\[
 \operatorname{In}(\phi)
   =\{r\in R:|\phi^{-1}(\{r\})|\geq\aleph_0\}.
\]
A run of \(\mathfrak A_0\) on \((v,T_x)\) is a map \(r:T_x\to S\) satisfying
\[
 (r(y0),r(y1))\in M(r(y),v(y))
\]
for every \(y\in T_x\).  Definition 1.5 accepts the valued tree exactly when
\[
 \exists r\left[
  r(x)\in S_0
  \ \land\
  \forall\pi\subseteq T_x\,
   \bigl(\pi\text{ is a path}
      \longrightarrow
      \operatorname{In}(r\mathbin{\upharpoonright}\pi)\in F\bigr)
 \right].
\]
Thus ordinary acceptance has the prefix \(\exists r\,\forall\pi\), and the
acceptance datum specified in Definition 1.5 is the family
\(F\subseteq\mathcal P(S)\).  The
finite pair sequence introduced later belongs to a dual construction; it is
not Definition 1.5's acceptance field.  The language \(T(\mathfrak A)\)
contains accepted valued subtrees rooted at arbitrary nodes \(x\), not only
full-root trees.  It is invariant under the evident re-rooting isomorphisms
\cite[pp.~5--7]{Rabin1969}.

Lemma 1.2 reduces a finite initial set to one initial state.  Theorem 1.3 gives
effective union and intersection.  Theorem 1.4 effectively projects away one
label coordinate and, conversely, adds a coordinate left unconstrained at
every node.  Theorem 1.5 gives effective complementation, and Theorem 1.6
effectively decides whether \(T(\mathfrak A)=\varnothing\).  The paper calls
the last two theorems difficult and proves them in Chapter III; the logical
decision argument depends on both
\cite[pp.~6--7]{Rabin1969}.

\subsubsection{Formula tracks and both definability directions}

For \(A\subseteq T\), let \(\chi_A:T\to\{0,1\}\) be its characteristic
function.  Definition 1.7 sends an \(n\)-tuple
\(\bar A=(A_1,\ldots,A_n)\) to the full-root valued tree
\[
 \tau(\bar A)=(v_{\bar A},T),
 \qquad
 \Gamma_n=\{0,1\}^n,
 \qquad
 v_{\bar A}(x)
   =\bigl(\chi_{A_1}(x),\ldots,\chi_{A_n}(x)\bigr).
\]
This is a bijection between \(\mathcal P(T)^n\) and the full-root
\(\Gamma_n\)-valued trees.  Theorem 1.7 effectively assigns to every
\(S2S\)-formula \(G(A_1,\ldots,A_n)\) an automaton
\(\mathfrak A_G\) such that
\[
 \tau(\bar A)\in T(\mathfrak A_G)
 \quad\Longleftrightarrow\quad
 \mathfrak N_2\models G(\bar A).
\]

The construction first eliminates equality and the two orders in favor of set
membership.  It then moves the set quantifiers outward and makes the remaining
matrix a Boolean combination of principal formulas
\[
 \exists x\,
 \bigwedge_{j=1}^{k} xw_j\,\eta_j\,A_{i_j},
 \qquad
 k\in\mathbb N_{>0},
 \quad
 w_j\in T,
 \quad
 i_j\in\{1,\ldots,n\},
 \quad
 \eta_j\in\{\in,\notin\},
\]
where \(xw_j\) is word concatenation, equivalently the indicated finite
successor iterate.  Section 3.10 later constructs an automaton for every such
principal formula; Corollary 3.13 explicitly fills this deferred
\enquote{small gap}.  Boolean combinations use union, intersection, and
complement.  Existential set quantification projects away one bit track;
universal set quantification uses complement and projection
\cite[pp.~7--9, 30--31]{Rabin1969}.

The source-specific finite witness used in the next example is fixed before
use.  A finite tree here is a finite set \(E\subseteq T\) with
\(\Lambda\in E\), closed under predecessors.  Put
\[
 \operatorname{Ft}(E)=\{x\in E:x0\notin E\ \land\ x1\notin E\}.
\]
Rabin calls this the terminal-node frontier and calls \(E\)
\emph{frontiered} when every \(x\in E\setminus\operatorname{Ft}(E)\) has both
children \(x0,x1\in E\).  Frontieredness is not redundant with predecessor
closure, which permits a node with exactly one child.  A finite valued tree is
a pair \((v,E)\) with
\(v:E\setminus\operatorname{Ft}(E)\to\Sigma\); the frontier is reserved for
the external run value defined with the finite-tree automaton below.  That
automaton and Lemma~3.12 use the pairs for which \(E\) is frontiered.  This
rooted frontiered source object, rather than an arbitrary finite subset of the
infinite tree, is what the lemma lifts.

For a concrete two-track instance, consider
\[
 P(A,B)=\exists x\,(x\in A\land r_0(x)\in B).
\]
The alphabet is \(\{0,1\}^2\).  The track tree satisfies \(P\) exactly when
some node has first bit \(1\) and its left child has second bit \(1\).  A
frontiered finite tree in Section 3.10's sense can be chosen so that both that
node and child are labelled, thereby recording the witness condition.  The
finite-tree definability fact used there, followed by Lemma 3.12, lifts this
finite witness language to an infinite-tree automaton.  This is a handbook-worked instance
of the source construction, not a transition table printed by Rabin.

The converse is Theorem 1.10.  After coding the finite state set into
\(\{0,1\}^m\), monadic tracks \(B_1,\ldots,B_m\) encode a state-valued run.
After Lemma 1.2's one-initial-state normalization, an \(S2S\)-formula says that
the root carries that state, every node and its two children obey the move
table, and, for every set \(C\) that is a path, the set of states occurring
infinitely often along \(C\) belongs to \(F\).
The resulting set-quantifier prefix is
\[
 \exists B_1\cdots\exists B_m\,\forall C,
\]
with a set-quantifier-free matrix.  It defines exactly the relation represented
by the automaton and yields Corollary 1.11 and its dual
\cite[pp.~9--10]{Rabin1969}.

For a closed sentence \(G\), Rabin first writes \(G\) without loss as
\(\exists A_1\,F(A_1)\).  Theorem 1.7 produces an automaton representing
\(F\), and \(G\) is true exactly when that automaton's language is nonempty,
which Theorem 1.6 decides.  Rabin's own proof has no additional regular-tree
basis step at this point
\cite[pp.~9--10]{Rabin1969}.

\subsubsection{Theorem 3.1 and its product construction}

If a path is enumerated by increasing length as
\(\pi=\{x_0,x_1,\ldots\}\), write
\((r\mathbin{\upharpoonright}\pi)_\omega(n)=r(x_n)\).  Theorem 3.1 begins
with a tree transition table \(\mathfrak A_0=\langle S,M\rangle\) and an
automaton-definable language \(B\subseteq S^\omega\), meaning a language
recognized by a finite \(\omega\)-word automaton.  It says that the
valued trees admitting one run \(r\) for which
\[
 \forall\pi\;
 (r\mathbin{\upharpoonright}\pi)_\omega\in B
\]
form an effectively tree-automaton-definable language.

The proof invokes McNaughton's external equivalence to choose the following
deterministic Muller word automaton for \(B\):
\[
 \mathfrak B=\langle U,K,u_0,H\rangle,
 \qquad
 \varnothing\neq U\text{ finite},
 \quad u_0\in U,
 \quad K:U\times S\longrightarrow U\text{ total},
 \qquad H\subseteq\mathcal P(U),
\]
For \(b=b_0b_1\cdots\in S^\omega\), its unique run
\(\rho_b:\mathbb N\to U\) satisfies
\[
 \rho_b(0)=u_0,
 \qquad
 \rho_b(n+1)=K(\rho_b(n),b_n),
 \qquad
 b\in B\Longleftrightarrow\operatorname{In}(\rho_b)\in H.
\]
The proof forms the product state set \(U\times S\), the initial set
\(\{u_0\}\times S\), and assigns at \(((u,s),a)\) the transition set
\[
 \bigl\{
      ((K(u,s),s_1),(K(u,s),s_2)):
      (s_1,s_2)\in M(s,a)
 \bigr\}.
\]
Along every branch the first coordinate is the unique word-automaton run on
the second-coordinate state sequence.  The designated family is
\[
 \{C\subseteq U\times S:p_1(C)\in H\},
 \qquad
 p_1(C)=\{u:\exists s\;(u,s)\in C\}.
\]
Determinism is what makes the first-coordinate histories shared coherently
before a branch splits.  The resulting product proves Theorem 3.1, which is
used later in the simultaneous-run constructions
\cite{McNaughton1966}\cite[pp.~17--19]{Rabin1969}.

\subsubsection{Dual acceptance and complementation}

Definition 3.5 fixes a finite sequence
\[
 \Omega=((L_i,U_i))_{1\leq i\leq k}
\]
of pairs of subsets of a state set \(S\).  For a map \(\phi:X\to S\) from any
set \(X\), define
\[
 \phi\in[\Omega]
 \quad\Longleftrightarrow\quad
 \exists i\in\{1,\ldots,k\}\;
 \bigl(
   \operatorname{In}(\phi)\cap U_i\neq\varnothing
   \ \land\
   \operatorname{In}(\phi)\cap L_i=\varnothing
 \bigr).
\]
The map is of type \(\Omega\)-empty when
\[
 \phi\in[\Omega,e]
 \quad\Longleftrightarrow\quad
 \phi\in[\Omega]
 \ \land\
 \exists j\in\{1,\ldots,k\}\;
   \operatorname{rng}(\phi)\cap L_j=\varnothing.
\]
Definition 3.6 says that a table \(\mathfrak A_0=\langle S,M\rangle\), a
sequence \(\Omega\), and a state \(s\) dually accept \((v,T_x)\) when
\[
 (v,T_x)\in D(\mathfrak A_0,\Omega,s)
 \quad\Longleftrightarrow\quad
 \forall r\left[
  \begin{array}{c}
   r\text{ is a run of }\mathfrak A_0\text{ on }(v,T_x)\\[-1mm]
   {}\land r(x)=s
  \end{array}
  \ \longrightarrow\
  \exists\pi\subseteq T_x\;
    r\mathbin{\upharpoonright}\pi\in[\Omega]
 \right].
\]
This \(\forall r\,\exists\pi\) prefix is dual to ordinary acceptance's
\(\exists r\,\forall\pi\).  Because the move table may have no run from
\(s\), the universal implication can hold vacuously.  Lemma 3.6 effectively
represents the complement of every ordinary automaton language as one such
dual language.  The pair sequence therefore belongs to complementation, not
to the ordinary automaton's acceptance component \(F\)
\cite[pp.~24--26]{Rabin1969}.

\paragraph{Handbook-owned zero-length completion.}
Lemma 3.6 forms its pair list by enumerating
\(\mathcal P(S)\setminus(F\cup\{\varnothing\})\), so that list can be empty.
For the empty sequence \(\langle\rangle\), no path map belongs to
\([\langle\rangle]\); hence
\[
 D(\mathfrak A_0,\langle\rangle,s)
 \quad\Longleftrightarrow\quad
 \text{there is no run of }\mathfrak A_0\text{ rooted in }s.
\]
The printed induction begins at length one and does not discharge this case.
To enter that positive-length induction, define the one-pair sequence
\(\Omega^{\dagger}=((S,S))\).  Every map from an infinite path into the
finite nonempty set \(S\) has a nonempty infinitely recurring set, so it
cannot both meet and avoid \(S\).  Therefore
\[
 [\Omega^{\dagger}]=\varnothing,
 \qquad
 D(\mathfrak A_0,\Omega^{\dagger},s)
   =D(\mathfrak A_0,\langle\rangle,s).
\]
This is not a new acceptance condition: it is an exact one-pair representation
of the otherwise omitted impossible path type, and it is used below whenever
Lemma 3.6 returns no pairs.

The proof that every dual language is ordinarily automaton definable has the
following source architecture.

\begin{enumerate}
\item For a finite marker set \(Q\) and
  \(H\subseteq T_x\times Q\), the tree marked at \(H\) labels a node \(y\) by
  \((v(y),\{q:(y,q)\in H\})\).  Lemma 3.2 proves that an iterated family of
  local marked-tree acceptance obligations is automaton definable.  Its proof
  runs all activated copies simultaneously, merges copies of the same local
  automaton that reach the same state, and uses Theorem 3.1 to check the
  remaining condition on every path.

\item For a nonempty set \(H\), Definition 3.2 calls a mapping
  \(h\mapsto H_h\subseteq H\) well founded exactly when no infinite sequence
  \(h_1,h_2,\ldots\) satisfies
  \(h_{m+1}\in H_{h_m}\) for every \(m\).  Lemma 3.3 characterizes this by an
  ordinal decomposition in which every member of \(H_h\) lies in a lower
  stratum than \(h\).  Lemma 3.4 adds this well-foundedness requirement to
  Lemma 3.2 and remains automaton definable.

\item Lemmas 3.7 and 3.8 translate the local disjunctions needed in the dual
  argument into new dual languages.  Lemma 3.7 preserves the length of the
  pair sequence; Lemma 3.8 produces a sequence of the same length whose last
  lower set is empty.

\item For
  \(\Omega=((L_i,U_i))_{1\leq i\leq k}\) with \(L_k=\varnothing\), Lemmas
  3.9--3.10 express dual acceptance by iterated local obligations involving
  the shorter sequence
  \(\Omega_{k-1}=((L_i,U_i))_{1\leq i<k}\).  Lemma 3.11 treats arbitrary
  \(\Omega\): every local obligation either finds an \(\Omega\)-empty path or
  defers to a strictly later marked obligation, and those deferrals must be
  well founded.

\item After the preceding zero-length normalization, Section 3.11 inducts on
  \(k\geq1\) in the paper's two
  \enquote{half-steps}.  The \(L_k=\varnothing\) case uses Lemma 3.10, the
  induction hypothesis for \(k-1\), and Lemma 3.2.  The arbitrary case uses
  Lemma 3.11, Lemma 3.8 to reach a last-empty-lower-set sequence, and Lemma
  3.4 to enforce well-foundedness.  Lemma 3.6 then yields Theorem 1.5.  The
  normalization, rather than a silent assumption that the pair list is
  nonempty, supplies the source's missing initial case.
\end{enumerate}

Finite-tree results support three parts of the article.  Recall that Section
3.10's finite tree is a finite predecessor-closed \(E\subseteq T\) with
\(\Lambda\in E\), and that its frontier and frontieredness satisfy
\[
 \operatorname{Ft}(E)=\{x\in E:x0\notin E\ \land\ x1\notin E\},
 \qquad
 x\in E\setminus\operatorname{Ft}(E)\Longrightarrow x0,x1\in E.
\]
Rabin calls \(\operatorname{Ft}(E)\) the terminal-node frontier and calls
\(E\) \emph{frontiered} exactly when the second displayed condition holds.
This adjective names a property of \(E\), not a second tree object.  It is not
redundant with predecessor closure, which permits a node with exactly one
child; the finite-tree automaton, lift, and emptiness base below require both
children at every node outside the frontier.
For the finite-tree automaton below, let \(E\) be frontiered.  Its finite
valuation has the previously fixed type
\(v:E\setminus\operatorname{Ft}(E)\to\Sigma\).  The source's
finite-tree automaton is
\[
 \langle S,M,s_0,f\rangle,
 \qquad f\notin S,
 \qquad
M:S\times\Sigma\longrightarrow
  \mathcal P\bigl((S\times S)\cup\{(f,f)\}\bigr).
\]
A run on this frontiered finite valued tree \((v,E)\) is a map
\[
 r:E\longrightarrow S\cup\{f\}
\]
such that \(r(x)=f\) for every \(x\in\operatorname{Ft}(E)\), while for every
\(x\in E\setminus\operatorname{Ft}(E)\),
\[
 r(x)\in S,
 \qquad
 (r(x0),r(x1))\in M(r(x),v(x)).
\]
The automaton accepts the frontiered input \((v,E)\) exactly when such a run
satisfies \(r(\Lambda)=s_0\).  Thus the distinguished external symbol
\(f\notin S\) is a frontier value, not an ordinary state, and the finite run is neither the
earlier \(S\)-valued infinite-tree run nor the host's bottom-up ranked-tree
run.  Because the printed transition codomain contains no mixed pair
\((f,s)\) or \((s,f)\) with \(s\in S\), a frontiered input having exactly one
frontier child at a node is a valid finite tree but has no run at that node; the
source's codomain is not silently enlarged.  For the formula translation,
Lemma 3.12 lifts any automaton-definable collection of
frontiered finite valued trees to the infinite valued trees having an accepted
frontiered finite initial subtree, and Corollary 3.13 uses that lift for the
principal formulas in Theorem 1.7.
For complementation, Lemma 3.14 uses K\"onig's Infinity Lemma: if every
infinite run meets a set of node--state obligations, then some finite subset
already meets every run.  Corollary 3.15 turns the resulting universal hitting
condition into an ordinary tree automaton, supplying the \(k=1\) base.
Section 3.12 then reuses finite witnesses in the emptiness decision described
below
\cite[pp.~30--32]{Rabin1969}.

Transfinite induction is therefore present at a precise locus.  Lemma 3.3
constructs the ordinal strata of a well-founded obligation mapping, and the
converse direction of Lemma 3.11 constructs its stages by transfinite
induction.  It is not an adequate proof outline to replace the marked trees,
simultaneous finite-state runs, dual languages, two half-steps, and induction
on \(k\) by the unqualified phrase
\enquote{a transfinite complementation proof}
\cite[pp.~19--32]{Rabin1969}.

\subsubsection{Rabin's actual emptiness argument}

Section 3.12 does not use a regular-tree basis theorem.  By Lemma 1.2 the
input automaton may be written
\(\mathfrak A=\langle S,M,\{s_0\},F\rangle\).  Over the one-letter alphabet
\(\{\bullet\}\), define
\[
 M^\bullet(s,\bullet)
   =\bigcup_{a\in\Sigma}M(s,a).
\]
Also define the valued full tree
\[
 v_\bullet:T\longrightarrow\{\bullet\},
 \qquad v_\bullet(x)=\bullet,
 \qquad \mathbf t_\bullet=(v_\bullet,T)\in V_{\{\bullet\}}.
\]
Every run of this table on \(\mathbf t_\bullet\) is a
run of \(\mathfrak A\) on some \(\Sigma\)-valuation, and conversely.  Hence
\[
 T(\mathfrak A)\neq\varnothing
 \quad\Longleftrightarrow\quad
 \mathbf t_\bullet\text{ is accepted by }
 \langle S,M^\bullet,\{s_0\},F\rangle.
\]
Theorem 1.5 constructs an automaton \(\mathfrak C\) for the complement of this
one-letter language.  Lemma 3.6 then constructs a table
\(\mathfrak B_0\) with finite state set \(R_0\), a sequence \(\Omega\) of pairs
of subsets of \(R_0\), and a state \(s'_0\in R_0\) such that
\[
 D(\mathfrak B_0,\Omega,s'_0)
   =V_{\{\bullet\}}-T(\mathfrak C),
\]
which is the original one-letter language.  Original emptiness is therefore
reduced to deciding whether \(\mathbf t_\bullet\) belongs to this specified dual
language.

That decision repeats the complementation proof's two half-steps, now as finite
state calculations.  In each recursive question below, let \(R\) be the state
set of the current finite one-letter table and let
\(\Omega=((L_i,U_i))_{1\leq i\leq k}\) be its current sequence of pairs of
subsets of \(R\); the state set can change under the effective reductions.  If
Lemma 3.6 produced \(k=0\), first replace its empty sequence by
\(((R,R))\) using the exact equivalence above.  Thus the following recursion
starts with \(k\geq1\).
For a current table \(\langle R,N\rangle\) and \(R'\subseteq R\), Definition
3.7's restricted table retains \(R\) as its formal state set and replaces
\(N(s)\) by
\[
 N_{R'}(s)=N(s)\cap(R'\times R').
\]
Its runs are exactly the runs of the original table whose states at nodes
strictly below the root lie in \(R'\).  This distinction matters when the root
state itself lies outside \(R'\).

The base case needs finite restrictions of the current one-letter table, not
the external-\(f\) finite-tree-automaton runs above.  For a finite frontiered
tree \(E\), a \emph{finite table-run rooted at \(s\in R\)} is a map
\[
 r:E\longrightarrow R
\]
such that \(r(\Lambda)=s\) and
\[
 (r(x0),r(x1))\in N(r(x))
 \qquad\bigl(x\in E\setminus\operatorname{Ft}(E)\bigr).
\]
Frontier nodes carry ordinary states and impose no outgoing-transition
condition.  This object is a finite restriction of a run of the current table;
it is neither an infinite run on \(T_x\) nor a run of
\(\langle S,M,s_0,f\rangle\).  The distinction is required in the terminal
decision clause below.

\begin{enumerate}
\item For arbitrary length-\(k\) \(\Omega\), the transfinite family from Lemma
  3.11 collapses over a finite state set to disjoint strata
  \(H_0,H_1,\ldots\); after at most \(|R|\) nonempty strata, no state remains
  to be added.  Lemma 3.8 reduces
  each membership question needed to compute a stage to one with a
  length-\(k\) sequence whose last lower set is empty.

\item If \(L_k=\varnothing\) and \(k>1\), the exhaustive Lemma 3.9 reduction
  ranges over all \(H\subseteq U_k\).  For \(H=\varnothing\), the clause is
  the direct shorter question
  \(\mathbf t_\bullet\in
    D(\langle R,N\rangle,\Omega_{k-1},s_0)\).  For each nonempty
  \(H\), Definition 3.7 with \(R'=R\setminus H\) converts the universal
  clauses for the states in \(H\cup\{s_0\}\) to dual questions for the
  length-\(k-1\) sequence \(\Omega_{k-1}\).  Rabin's display (3.5) prints
  only the nonempty choices, although its following reduction ranges over
  \(H\in\mathcal P(R)\); admitting \(H=\varnothing\) is a handbook-owned
  repair.  It covers both the no-run case and the nonvacuous case in which
  every run already supplies an \(\Omega_{k-1}\)-type path.

\item In the base case \(k=1\) and \(L_1=\varnothing\), Lemmas 3.9 and 3.14
  give an exact finite witness after a second handbook-owned edge repair.  Let
  \(m\geq0\), allow
  \(U_0=\{s_1,\ldots,s_m\}\subseteq U_1\) to be empty, and take finite
  frontiered trees
  \(E_0,\ldots,E_m\).  Here \(s_0\) is the state in the current base-case
  dual question.  For every \(0\leq i\leq m\), every finite table-run \(r\) on
  \(E_i\) rooted at \(s_i\) must satisfy
  \[
   r(E_i\setminus\{\Lambda\})\cap U_0\neq\varnothing.
  \]
  Existence
  of each
  such finite tree is expressible in the weak monadic second-order theory of
  two successors.  Rabin cites Doner and Thatcher--Wright for the decidability
  of that finite-tree theory.  When \(m=0\), the condition says exactly that
  the sole finite tree \(E_0\) admits no finite table-run rooted at \(s_0\).
  K\"onig's Infinity
  Lemma supplies such a finite obstruction exactly when no infinite run exists.
  Rabin's p.~34 instead prints \enquote{nonempty}; that excludes dual
  acceptance's permitted vacuous branch and is not retained as an exhaustive
  equivalence.
\end{enumerate}

With these two explicit edge repairs and the zero-length normalization, the
source construction closes Theorem 1.6 effectively.  Rabin's printed
elementary-recursive
qualification records the article's historical claim; the handbook does not
retain it, for the later lower-bound reason stated above.  The
finite-tree/weak-theory result is an external theorem at the base of Rabin's
source proof; it does not make the input to the \(S2S\) automaton finite and
does not turn full monadic quantification into weak quantification
\cite[pp.~32--34]{Rabin1969}.

\subsubsection{Applications, later route, and proof boundary}

Chapter II applies the binary-tree interpretation well beyond the main
sentence decision.  Theorem 2.1 treats the monadic second-order theory of the
class of countable linear orders; Corollary 2.2 treats the weak theory of all
linear orders, and Corollary 2.3 the monadic second-order theory of the class
of countable well-orders.  Theorem 2.4 treats the class of unary-function
structures with countable domain, with Corollary 2.5 giving the weak theory of
all unary-function structures.  Theorems 2.6--2.9 code closed and
\(F_\sigma\) subsets of Cantor space and derive decision results for ordered
Cantor space with those set variables, the lattice of \(F_\sigma\) sets with
the closed sets distinguished, and the analogous language over the unit
interval.  Theorems 2.10--2.11 use Stone representation for the theory of
countable Boolean algebras with quantification over ideals and the
first-order theory of arbitrary Boolean algebras with a sequence of
distinguished ideals.  Section 2.5 codes a Gale--Stewart strategy by a node set
and expresses in \(S2S\) the determinacy statement for \(F_\sigma\) payoff
sets; Rabin presents this as placing Wolfe's theorem among statements decided
by the procedure
\cite[pp.~10--17]{Rabin1969}.

The presentation above reconstructs Rabin's structures, both logical
translations, one worked track encoding, the path-language product, the exact
ordinary/dual acceptance distinction, the marked-tree and well-founded
complement architecture, the two half-steps, the precise transfinite locus,
the finite-tree lemmas, the source's emptiness reduction, and the application
boundary.  It does not reproduce every transition table and every induction
inside the article's fourteen-page complementation development; those
constructions remain in the cited complete source.

Thomas's report says expressly that its Sections~6.1--6.2 give a later,
game-theoretic and fairly self-contained proof assembled from post-1969 work;
it is not presented as Rabin's original proof.  Let \(Q\) be a finite nonempty
state set, \(\Sigma\) a finite alphabet, \(q_0\in Q\), and
\(\mathcal F\subseteq\mathcal P(Q)\).  Put \(T=\{0,1\}^*\).  In
Definition~6.1 a Muller tree automaton is
\[
 \mathcal A=(Q,\Sigma,q_0,\Delta,\mathcal F),
 \qquad
 \Delta\subseteq Q\times\Sigma\times Q\times Q.
\]
For an input \(t:T\to\Sigma\), a run is a map \(r:T\to Q\) with
\[
 r(\epsilon)=q_0,
 \qquad
 (r(u),t(u),r(u0),r(u1))\in\Delta
 \quad(u\in T).
\]
For a direction stream \(\pi\in\{0,1\}^\omega\), let
\(\pi\mathbin{\upharpoonright}0=\epsilon\), let
\(\pi\mathbin{\upharpoonright}n\in\{0,1\}^*\) be its length-\(n\) prefix, and
define the induced state stream
\[
 r_\pi:\mathbb N\longrightarrow Q,
 \qquad
 r_\pi(n)=r(\pi\mathbin{\upharpoonright}n).
\]
Thomas's \(\pi\) is a direction stream; it is not Rabin's earlier node-set
path.  The run is successful exactly when
\[
 \forall\pi\in\{0,1\}^\omega\;
   \operatorname{In}(r_\pi)\in\mathcal F,
\]
and \(\mathcal A\) accepts \(t\) when such a run exists.  Thus tree acceptance
has the essential prefix \(\exists r\,\forall\pi\): one global run must satisfy
the word condition on every path.  Examples~6.2--6.3 exhibit complementary
languages that separate B\"uchi from Muller tree recognition
\cite[pp.~43--45]{Thomas1996Report}.

Definition~6.4 replaces the arbitrary Muller family by a finite strict Rabin
chain.  Write its length as \(m\geq0\).  For \(m>0\) it is
\[
 E_1\subsetneq F_1\subsetneq\cdots
   \subsetneq E_m\subsetneq F_m\subseteq Q.
\]
A path accepts exactly when, for some \(1\leq i\leq m\),
\[
 \operatorname{In}(r_\pi)\cap E_i=\varnothing
 \quad\text{and}\quad
 \operatorname{In}(r_\pi)\cap F_i\neq\varnothing.
\]
For \(m=0\), the chain is empty and the existential condition over \(i\) is
false, so every path rejects.  The report does not separately state this
endpoint, but Theorem~6.5 instructs the reader to drop equal adjacent pairs and
can thereby delete every pair; the empty-chain clause is the required
handbook-owned completion of that construction.  For \(m>0\), the source also
calls the condition a parity condition and reads the successive difference
sets by odd and even indices.  Theorem~6.5 converts a Muller tree automaton to
this form, possibly with the empty chain.  Its state is an \emph{order-vector with hit}: a
permutation of \(Q\), ordered by last appearance, together with the former
position of the newly visited state.  This object is finite memory added to an
automaton state, not an input tree and not a path.  Remark~6.6 recovers the
infinitely recurring state set from the eventual hit positions, which is the
specific reason the added memory implements the Muller condition
\cite[pp.~45--47]{Thomas1996Report}.

Before the game reduction, one premise omitted by the report must be supplied.
Definition~6.1 permits \(\Delta\) to be partial, whereas the game theorem below
requires a nonempty outgoing-edge set at every position.  For \(q\in Q\) and
\(a\in\Sigma\), put
\[
 \Delta(q,a)=\{(p_0,p_1)\in Q^2:
                  (q,a,p_0,p_1)\in\Delta\}.
\]
Choose \(\bot\notin Q\), let \(Q^\bot=Q\cup\{\bot\}\), retain the same chain
sets \(E_i,F_i\) as subsets of \(Q^\bot\), and set
\[
 \Delta^\bot
 =\Delta
  \cup\{(q,a,\bot,\bot):q\in Q,\ a\in\Sigma,
                            \ \Delta(q,a)=\varnothing\}
  \cup\{(\bot,a,\bot,\bot):a\in\Sigma\}.
\]
Write \(\mathcal A^\bot\) for this completed chain tree automaton.  The notation
names a chain tree automaton, not a game or a run; it is needed only to place
the later game inside its stated hypothesis, differs from \(\mathcal A\) only
where a state-letter pair had no transition, and is used throughout the game
argument below.  Every run of \(\mathcal A\) remains a run of
\(\mathcal A^\bot\).  Conversely, if a run uses an added transition, every path
through either new child is thereafter constantly \(\bot\).  Since
\(\bot\notin F_i\) for every \(i\), such a path satisfies no chain pair, so a
successful run cannot use an added transition.  Hence the recognized tree
language and the chain, possibly empty, are unchanged.  Consequently every position in
the game constructed below has a nonempty outgoing-move set.  This is a
handbook-owned completion of the source's missing game premise
\cite[pp.~44, 50]{Thomas1996Report}.

For the remainder of the game argument, replace \(\mathcal A\) by the
language-equivalent \(\mathcal A^\bot\) and again write \(Q\), \(\mathcal A\),
and \(\Delta\) for the completed state set \(Q^\bot\), automaton
\(\mathcal A^\bot\), and transition relation \(\Delta^\bot\), respectively.
For this completed chain automaton and input tree \(t\), Section~6.2 forms the
game \(\Gamma_{\mathcal A,t}\).  Its Automaton positions are the pairs
\((u,q)\in T\times Q\).  From \((u,q)\), Automaton chooses a transition
\[
 \tau=(q,t(u),p_0,p_1)\in\Delta
\]
and moves to the Pathfinder position \((u,\tau)\).  From there Pathfinder
chooses \(d\in\{0,1\}\), and play continues at \((ud,p_d)\).  The initial
position is \((\epsilon,q_0)\).  Completion makes every play infinite.
Automaton wins exactly when the generated state stream satisfies the chain
condition above; Pathfinder wins otherwise.  Remark~6.7 proves that
\(\mathcal A\) accepts \(t\) exactly
when Automaton has a winning strategy from the initial position.  A strategy
is \emph{memoryless} when its choice depends only on the current game position,
not on the earlier play.  Theorem~6.8 says that one player has such a winning
strategy from every position.  In Theorem~6.9, failure of acceptance therefore
gives Pathfinder a memoryless winning strategy.  Each of Pathfinder's local
instructions is a function
\[
 i:\Sigma\times\Delta\longrightarrow\{0,1\},
\]
where at a node \(u\) only the compatible pair \((t(u),\tau)\), with
\(\tau=(q,t(u),p_0,p_1)\), is queried.  A tree labelled by these finitely many
functions therefore encodes the strategy at all nodes.  Closure under negation
and projection produces a sequential
deterministic Muller automaton that universally checks the transition sequences
consistent with the instructions along one path.  A deterministic Muller tree
automaton runs that word automaton on every path, and nondeterminism guesses the
instruction labels.  The resulting complement automaton depends
only on \(\mathcal A\), not on the particular input \(t\).  Since completion is
effective and language-preserving, this automaton complements the original
possibly partial automaton as well.  The completed game has nonempty finite
outdegree at every position, so it lies within the hypothesis of
the game theorem.  If the chain is empty, its path condition is false and
Pathfinder wins every completed play; this is the trivial base.  For a
nonempty chain, Theorem~6.13 proves the needed determinacy by induction on its
entries: it collects the positions from which one player can force a finite
visit to a target set and recurses on the induced complementary subgame.
Theorem~6.16 makes the construction effective on finite game graphs
\cite[pp.~47--55]{Thomas1996Report}.

Definition~6.17 calls a \(\Sigma\)-labelled tree \emph{regular} exactly when a
deterministic finite word automaton over directions \(\{0,1\}\), equipped with
an output map to \(\Sigma\), generates its node labels.  Concretely, this means
there are a finite nonempty set \(P\), an initial state \(p_0\in P\), a total
map \(\delta:P\times\{0,1\}\to P\), and an output map
\(\lambda:P\to\Sigma\).  Extend \(\delta\) to direction words by
\[
 \delta^*(p,\epsilon)=p,
 \qquad
 \delta^*(p,ud)=\delta(\delta^*(p,u),d)
 \quad(u\in\{0,1\}^*,\ d\in\{0,1\}),
\]
and require
\[
 t(u)=\lambda(\delta^*(p_0,u))\qquad(u\in T).
\]
Regularity is a property of
the input tree, not an acceptance mode.  Theorem~6.18 first removes the input
by letting an automaton guess labels.  The handbook completes that automaton as
above before the resulting deadlock-free finite game is solved.  An empty chain
recognizes the empty language.  Otherwise, when the language is nonempty, an
effective memoryless winning strategy fixes
one transition at every automaton position and therefore supplies a finite
deterministic generator for an accepted regular tree.  Hence emptiness is
decidable and every nonempty recognized language has a regular member
\cite[pp.~55--57]{Thomas1996Report}.

Finally, for \(i\in\{0,1\}\), define the successor function
\[
 S_i:T\longrightarrow T,
 \qquad S_i(u)=ui.
\]
These structure functions are unrelated to Rabin's earlier initial-state set
\(S_0\subseteq S\).  Section~6.3 works in
\(\mathfrak T=(T,S_0,S_1)\).  Theorem~6.19 constructs a Muller tree
automaton from an MSO formula by induction in the singleton-set presentation:
atomic formulae supply the base automata, disjunction and existential set
quantification use nondeterminism, and negation uses Theorem~6.9 after the
effective language-preserving completion above.  Conversely,
existential monadic tracks encode a run and a universal path set enforces its
acceptance.  For a sentence the first
construction is input-free; after the same completion, Theorem~6.18 decides
whether it has a successful run.  Theorem~6.20 therefore concludes that
\(S2S\) is decidable.
The regular-tree conclusion also supplies the report's regular-model
consequence for translated modal and branching-time logics
\cite[pp.~57--59]{Thomas1996Report}.

This is Thomas's later complementation-and-basis route.  It is not Rabin's
Section~3.12 emptiness proof, is not projected backward onto the 1969 article,
and does not replace the Rabin reconstruction above.  The registered 1997
chapter is not part of the source basis used here, so every locator in this
Thomas presentation belongs to the May 1996 report precursor, not to the final
chapter.
